\section{Initial symbols satisfy the incoming filtration}
\label{sec:initial-membership}

\paragraph{IN1--IN2: the exact statement.}
Assume the archive's \texttt{DataII} hypotheses with $C\geq1$, and put
$\mathbb K=\operatorname{Frac}(O)$, $k=r+K$.  For every actual symbol
$S=\mathrm{SPR}\,B_b$, every physical digit block, every selected normal
joint projector $e_\kappa$, and every normal coefficient, we prove
\[
 H=[Z^\nu](1\otimes e_\kappa)\operatorname{lev}_{\rho,d_P,d_N}(S)
 \quad\Longrightarrow\quad H\in X_D(0,1/8).
 \tag{IN1}
\]
This is precisely \texttt{Lemma5aAt}, the initial half of incoming
\texttt{MemAt}.  Here $v_T$ denotes the minimum exponent of a
generic-conjugated rank-one torus acting on the signed letters in $T$;
$v_T(0)=+\infty$.  The required conditions are plane multidegree $d_P$,
chain bounds $\theta_{j,i}$ after every frame-digit extraction, and link
bounds $\theta_{j,u}-s_{j,u}/8$.

Expand each segment factor into its defining plane--normal generators
$A_jG_j$, and expand $\mathrm{SPR}$ into monomials $X^\gamma$.
Each $\gamma$ uses only plane coordinates at even, hence vector, letters:
every determinant defining $\mathrm{SPR}$ has this property.
Retain the actual total degree
$d^{\rm tot}=\mathrm{G2d}+\deg\gamma$ through the compatibility condition
\[
 \sum_c\rho(x,c)+q_x\bigl(d_P(x)+d_N(x)\bigr)
       =d^{\rm tot}(\operatorname{pos}x),
 \qquad q_x=p^{z_j-\operatorname{pos}x}.
 \tag{IN2}
\]
Individual expanded summands need not have that degree.  We first prove
membership only on blocks satisfying (IN2), extend by linearity, and
then use the actual symbol's homogeneity: all incompatible blocks vanish.

\paragraph{IN3--IN4: the normal projector supplies the threshold.}
For one product summand, put $b=\lfloor(l+\gamma)_{\rm plane}/q\rfloor$.
Its digit block is zero unless $\rho\equiv l+\gamma\pmod q$; otherwise it is
$A_*\otimes G_*$.  The normal factor is the prescribed homogeneous
component of the product of renamed $G_j$'s.  Normal labels disappear
because their exponents are $<q$, and $G_j$ has no preceding-frame variables.
Consequently
\[
 [Z^\nu](1\otimes e_\kappa)(A_*\otimes G_*)
             =[Z^\nu](e_\kappa G_*)\,A_*.
 \tag{IN3}
\]
If $e_\kappa G_*\neq0$, every selected prefix type is present.
Indeed canonical isotypic projectors commute for nested or disjoint
letter sets, so their product is the unique joint projector.
The generic maximum weight of a nonzero highest-weight type $\lambda$
is $\lambda_1$; homogeneous extraction cannot increase it.
Thus the original normal-prefix upper bounds also bound the corresponding
$\kappa_{j,u,1}$.

Write $x_{j,u}$ for the ordered nonpreceding-frame letters,
$B_{j,u}=\mathrm{Bo}(j,x_{j,u})$, and
$s_{j,u}=\lfloor B_{j,u}/(Yk^2)\rfloor$.  The defining pair inequalities give
\[
 v_{<x_{j,u}}(A_j)\geq\theta_{j,u},\qquad
 \theta_{j,u}=\max\left\{
 \begin{array}{l}
 B_{j,u}/2,\\
 \kappa_{j,u+1,1}+m_*+s_{j,u}\quad\text{if }x_{j,u}\text{ is even},\\
 \kappa_{j,u-1,1}+m_*+s_{j,u}\quad\text{if }u>0,
                         \ x_{j,u-1}\text{ is odd}
 \end{array}\right\}.
 \tag{IN4}
\]
Absent terms are omitted.  For the last line, the supported sets
$\{x\leq x_{j,u-1}\}$ and $\{x<x_{j,u}\}$ coincide, and an odd
predecessor leaves the oldest preceding vector, hence $B$ and $s$, unchanged.

\paragraph{IN5--IN6: frame digits cost no chain order.}
Give each segment independent intermediate letters and current-frame
labels, and give its preceding frame independent variables $w_{j-1,m}$.
Let $\chi_j$ rename $A_j$ into these variables and
$\Pi=\prod_j\chi_jA_j$.  The substitution $\phi$ fixes intermediate letters
and sends
\[
 \mathrm{lab}_{j,m,a}\mapsto X_{j,m,a},\qquad
 w_{j,m,a}\mapsto X_{j,m,a}^{P_j},\qquad
 P_j=p^{z_{j+1}-z_j}.
\]
Exactly $A_*=[d_P](X^b\phi\Pi)=X^b\phi\Pi'$, where $\Pi'$ keeps
the model multidegrees mapping to $d_P$.  Such filters commute with
the generic torus.  On the model chain consisting of $w_{j-1,m}$ and
the first $i$ intermediate letters, only $\chi_jA_j$ contributes; therefore
$v_{C_{j,i}}(\Pi')\geq\theta_{j,i}$.

Expand $\Pi'=\sum_M\mathrm{lab}^M Q_M$.  Labels lie outside every chain,
so each $Q_M$ retains its chain bound and uses only intermediate and $w$
variables.  Direct division of monomial exponents gives
\[
 \operatorname{Fdig}_{\sigma}(X^\beta\phi Q_M)=
 \begin{cases}
 X^{\operatorname{lift}\lfloor\beta/P\rfloor}Q_M,
                   &\beta\bmod P=\sigma,\\
 0,&\text{otherwise}.
 \end{cases}
 \tag{IN5}
\]
Here $P=1$ on intermediate letters.  With $\beta=b+\phi\!\operatorname{exp}(M)$,
the extra monomial has no covector variables inside the chain: labels
are frame vectors, and $b$ has no odd intermediate exponents.
Its order is nonnegative.  Summation proves the exact chain bound.

\paragraph{IN7--IN8: the sole loss is a preceding-frame degree.}
For a link ending at $x_{j,u}$, where $u=n_j+r-1-m$ and $m+2\leq r$,
remove the preceding frame $W_j$ from the original prefix.  If a
homogeneous piece has total $W_j$-degree $d_W$, this lowers order by
at most $d_W$: apply the inverse vector twist on $W_j$, whose Laurent
exponents are at least $-d_W$, then extract its coefficients.
Remaining letters rename into the link; all other factors contribute
only vector variables there.  A piece surviving $[d_P]$ satisfies
$P_{j-1}\deg_{W_{j,m'}}\leq d_P(j-1,m')-\deg_{j-1,m'}b$.  Hence
\[
 v_{T_{j,m}}(A_*)\geq\theta_{j,u}-\Delta_j,\qquad
 \Delta_j=\sum_{m'=0}^{r-1}
 \frac{\bigl(d_P(j-1,m')-\deg_{j-1,m'}b\bigr)_+}{p^{g_j}},
 \quad g_j=z_j-z_{j-1}.
 \tag{IN6}
\]
The fixed-degree relation (IN2) gives, for distinct integers
$t_{m'}=z_{j-1}-\operatorname{pos}(j-1,m')\geq0$,
$p^{t_{m'}}(d_P-\deg b)_+\leq\mathrm{G2d}\leq2a$.
Thus $\Delta_j\leq3a/p^{g_j}$, since $p\geq3$.
If $D=z_j-\operatorname{oldest}(j,x_{j,u})$, frame separation gives
$D+L_1+1\leq g_j$.  The source size hypotheses with $C\geq1$ give
\[
 a\leq B_{j,u}p^D,\qquad 48Yk^2\leq p^{L_1+1},\qquad
 1\leq B_{j,u}/(Yk^2)\leq2s_{j,u}.
\]
Consequently
\[
 \Delta_j\leq\frac{3a}{p^{g_j}}
 \leq\frac{3B_{j,u}}{p^{L_1+1}}
 \leq\frac{B_{j,u}}{16Yk^2}\leq\frac{s_{j,u}}8.
 \tag{IN7}
\]

\paragraph{IN9: assembly.}
Equations (IN3)--(IN7), plane homogeneity, and linearity prove (IN1).
This closes initial \texttt{MemAtG(0,1/8)} for the actual symbol family.
The following word-transport argument spends the separate operator slack.
